% This file is a chapter fragment.  It is included by version/detailed.tex.
%
% Status convention: every \lean{} field in this chapter is deliberately left
% empty and no leanok marker is set.  The chapter is a plan, written before the Lean
% declarations exist; a separate integration step fills the names in once the
% implementation pull requests have landed.  Until then the labels below are
% the stable citation points: they are referred to from the Lean sources and
% from the implementation issues, so they must not be renamed.

\chapter{The open--closed exact sequence}\label{chap:openclosed}

Let $X$ be a scheme, let $j\colon U\to X$ be an open immersion and let
$i\colon Z\to X$ be a closed immersion whose range is the complement of the
range of $j$.  The goal of this chapter is the short exact sequence
\[
  0\longrightarrow j_!j^*\mathcal F\longrightarrow\mathcal F
   \longrightarrow i_*i^*\mathcal F\longrightarrow0
\]
of sheaves on the small \'etale site of $X$, for an arbitrary sheaf
$\mathcal F$ with values in a Grothendieck abelian category of coefficients.
It is the statement \ref{thm:openclosed-exact} at the end of the chapter, and
everything before it is a step of its proof.

\section{Setting}\label{sec:openclosed-setting}

Throughout the chapter we fix the following data.

\begin{itemize}
\item A coefficient category $\mathcal A$, abelian and Grothendieck abelian.
      For the two conservativity statements
      \ref{lemma:openclosed-point-ulift} and
      \ref{prop:openclosed-conservative} we ask in addition that $\mathcal A$
      be concrete in the sense recalled in
      \ref{rmk:openclosed-universes}; the case of interest,
      $\mathcal A=\mathrm{Mod}_\Lambda$ for the coefficient ring $\Lambda$ of
      the $\ell$-adic formalism, satisfies this.
\item A scheme $X$.  For a scheme $Y$ we write $\et Y$ for its small \'etale
      site: the objects are the schemes \'etale over $Y$, the morphisms are
      the $Y$-morphisms, and the topology is generated by the jointly
      surjective families.  We write $\shv{\et Y}$ for the category of
      $\mathcal A$-valued sheaves on it.
\item A morphism of schemes $g\colon Y'\to Y$ induces an adjunction
      $g^*\dashv g_*$ between $\shv{\et Y}$ and $\shv{\et{Y'}}$, with
      \[
        (g_*\mathcal G)(V)=\mathcal G(V\times_YY')
        \qquad(V\in\et Y).
      \]
      Here $g_*$ is left exact and $g^*$ is exact.  If $g$ is an open
      immersion then $g^*$ has in addition a left adjoint $g_!$, the
      extension by zero, which is exact, and the unit
      $\mathrm{id}\to g^*g_!$ is an isomorphism.
\item An open immersion $j\colon U\to X$ and a closed immersion
      $i\colon Z\to X$ such that
      \[
        \operatorname{range}(i)=X\smallsetminus\operatorname{range}(j).
      \]
      We call such a pair $(j,i)$ a \emph{complementary pair} on $X$.  Only
      the disjointness of the two ranges is used before
      \ref{lemma:openclosed-point-factors}; the covering half of the
      condition enters exactly once, in
      \ref{lemma:openclosed-point-factors}.
\end{itemize}

For a sheaf $\mathcal F$ on $\et X$ we write
\[
  \alpha_{\mathcal F}\colon j_!j^*\mathcal F\to\mathcal F,
  \qquad
  \beta_{\mathcal F}\colon\mathcal F\to i_*i^*\mathcal F
\]
for the counit of $j_!\dashv j^*$ and the unit of $i^*\dashv i_*$.  These are
the two maps of \ref{thm:openclosed-exact}; no other map is constructed in
this chapter.

\begin{remark}[Universes]\label{rmk:openclosed-universes}
Schemes live in a universe $u$ and the coefficient category $\mathcal A$ is
large relative to it: $\mathcal A$ has morphisms in $\mathrm{Type}\,(u+1)$
and is concrete over $\mathrm{Type}\,(u+1)$, so that the forgetful functor
reflects isomorphisms and preserves filtered colimits at that universe.  The
geometric points of $\et X$, on the other hand, have fibres in
$\mathrm{Type}\,u$.  The two do not meet, and the standard conservativity
statement for a family of points is stated at a single universe.
\ref{lemma:openclosed-point-ulift} is what repairs the mismatch; every
stalk argument in this chapter is understood to take place after that
repair.  The extra concreteness hypotheses on $\mathcal A$ are therefore
carried by \ref{prop:openclosed-lowerStar-fully-faithful},
\ref{prop:openclosed-conservative} and \ref{thm:openclosed-exact}.  This is
expected rather than a defect: joint conservativity is a statement about
stalks, and stalks are sets.
\end{remark}

\section{Vanishing off the complement}\label{sec:openclosed-vanishing}

The first three statements say that a sheaf pushed forward from $Z$ has no
sections over $U$, and that everything formal follows from this.

\begin{lemma}[Disjoint ranges give an empty fibre product]
\label{lemma:openclosed-empty-pullback}
\uses{}
\lean{}
Let $f\colon V\to X$ and $g\colon W\to X$ be morphisms of schemes whose
ranges are disjoint.  Then the scheme $V\times_XW$ is empty.
\end{lemma}
\begin{proof}
A point $t$ of $V\times_XW$ has images $v\in V$ and $w\in W$ with
$f(v)=g(w)$, because the two projections are morphisms over $X$.  That point
of $X$ lies in both ranges, contradicting disjointness.
\end{proof}

\begin{lemma}[$j^*i_*$ vanishes]
\label{lemma:openclosed-vanishing-lowerStar}
\uses{lemma:openclosed-empty-pullback}
\lean{}
Let $j\colon U\to X$ and $i\colon Z\to X$ have disjoint ranges, with $j$ an
open immersion.  Then $j^*i_*\mathcal H=0$ for every
$\mathcal H\in\shv{\et Z}$.
\end{lemma}
\begin{proof}
For an open immersion $j$ the functor $j^*$ is restriction along the
base-change functor $\et X\to\et U$, so for $T\in\et U$, regarded as an
object of $\et X$ through $j$,
\[
  (j^*i_*\mathcal H)(T)=(i_*\mathcal H)(T)=\mathcal H(T\times_XZ).
\]
The structure morphism $T\to X$ factors through $j$, so its range is
contained in that of $j$ and hence disjoint from that of $i$; by
\ref{lemma:openclosed-empty-pullback} the scheme $T\times_XZ$ is empty.  The
sheaf condition for the empty covering of the empty scheme says that
$\mathcal H(\emptyset)$ is a terminal object, and in an additive category a
terminal object is a zero object.  A sheaf all of whose sections are zero is
a zero object of $\shv{\et U}$.
\end{proof}

\begin{lemma}[$i^*j_!$ vanishes]
\label{lemma:openclosed-vanishing-lowerShriek}
\uses{lemma:openclosed-vanishing-lowerStar}
\lean{}
In the situation of \ref{lemma:openclosed-vanishing-lowerStar},
$i^*j_!\mathcal G=0$ for every $\mathcal G\in\shv{\et U}$.
\end{lemma}
\begin{proof}
The point of the argument is that no section of $j_!\mathcal G$ is computed:
$j_!$ is a sheafification and its sections are not accessible.  Instead one
uses only adjunctions.  For every $\mathcal H\in\shv{\et Z}$,
\[
  \Hom(i^*j_!\mathcal G,\mathcal H)
  \cong\Hom(j_!\mathcal G,i_*\mathcal H)
  \cong\Hom(\mathcal G,j^*i_*\mathcal H)=0
\]
by \ref{lemma:openclosed-vanishing-lowerStar}.  Taking
$\mathcal H=i^*j_!\mathcal G$ shows that the identity of
$i^*j_!\mathcal G$ is zero, that is, $i^*j_!\mathcal G=0$.
\end{proof}

\section{Geometric points}\label{sec:openclosed-points}

\begin{lemma}[A point in the closed part factors through it]
\label{lemma:openclosed-point-factors}
\uses{}
\lean{}
Let $i\colon Z\to X$ be a closed immersion, let $K$ be a field and let
$x\colon\Spec K\to X$ be a morphism whose image point lies in the range of
$i$.  Then there is a unique $\tilde x\colon\Spec K\to Z$ with
$i\circ\tilde x=x$.

Consequently, if $(j,i)$ is a complementary pair on $X$, then every
$x\colon\Spec K\to X$ factors, uniquely, through exactly one of $j$ and $i$:
the geometric points of $X$ are the geometric points of $U$ together with the
geometric points of $Z$.
\end{lemma}
\begin{proof}
Uniqueness holds because a closed immersion is a monomorphism.  For
existence, let $p\in X$ be the image point and choose an affine open
$\Spec R\subseteq X$ containing $p$; the closed subscheme $Z$ meets it in
$\Spec R/I$ for an ideal $I\subseteq R$.  The morphism $x$ factors through
$\Spec R$ and corresponds to a ring map $R\to K$ whose kernel is the prime
$p$.  Since $p$ lies in the range of $i$ we have $I\subseteq p$, so
$R\to K$ kills $I$ and factors through $R/I$.  The resulting morphism
$\Spec K\to\Spec R/I=Z\cap\Spec R\subseteq Z$ is the required $\tilde x$; no
gluing is needed, because $\Spec K$ has a single point and therefore stays
inside the chosen chart.

For the second statement, the image point of $x$ lies either in the range of
$j$, in which case $x$ factors uniquely through the open immersion $j$ by the
universal property of an open subscheme, or in its complement, which is the
range of $i$.  The two cases are exclusive because the ranges are disjoint.
\end{proof}

\begin{lemma}[Raising a point a universe]
\label{lemma:openclosed-point-ulift}
\uses{}
\lean{}
Let $(\mathcal C,J)$ be a site and let $p$ be a point of it: a stalk functor
$\shv{\mathcal C}\to\mathrm{Type}\,w$ that preserves finite limits and all
colimits.  Then composing the fibre functor with the universe-raising
equivalence $\mathrm{Type}\,w\to\mathrm{Type}\,\max(w,w')$ is again a point of
$(\mathcal C,J)$, and its stalks are the old ones up to that equivalence.

Applied to $\et X$: the raised family of geometric points of $\et X$ is a
conservative family of points at the universe of the coefficients.  Concretely,
a morphism of sheaves on $\et X$ with values in $\mathcal A$ is an
isomorphism as soon as it is an isomorphism on the stalk at every geometric
point of $X$.
\end{lemma}
\begin{proof}
Universe raising is an equivalence of categories which preserves and reflects
all limits and colimits, so composing with it preserves every property in the
definition of a point.  The stalk functors of the raised point are the old
ones followed by an equivalence, hence detect exactly the same
isomorphisms.

For the second statement, conservativity of a family of points is checked
against a covering-sieve criterion: a family is conservative as soon as every
sieve $S$ on an object $V$ whose associated subobject of the sheaf represented
by $V$ has full stalk at every point of the family is a covering sieve.  That
criterion mentions no fibre universe, so the proof of
conservativity for the geometric points of $\et X$ at the universe of schemes
transcribes verbatim to the raised family.  What is used about $\et X$ is
that for a scheme $V$ \'etale over $X$ and a point $v\in V$ there is a
geometric point of $X$ through which $v$ is visible; this is the existence of
enough geometric points on the small \'etale site.
\end{proof}

\begin{lemma}[Stalks of a pullback]
\label{lemma:openclosed-stalk-upperStar}
\uses{}
\lean{}
Let $f\colon Y\to X$ be a morphism of schemes, let $\bar y$ be a geometric
point of $Y$ and let $\mathcal F\in\shv{\et X}$.  Then
\[
  (f^*\mathcal F)_{\bar y}\cong\mathcal F_{f\circ\bar y},
\]
naturally in $\mathcal F$.
\end{lemma}
\begin{proof}
The functor $f^*$ is the pullback of sheaves along the morphism of sites
given by base change $\et X\to\et Y$, $V\mapsto V\times_XY$.  Pulling a point
of $\et Y$ back along a morphism of sites gives a point of $\et X$, and the
stalk of a pulled-back sheaf at a point is by construction the stalk of the
sheaf at the transported point.  It remains to identify the transport of the
geometric point $\bar y$ along base change with the geometric point
$f\circ\bar y$ of $X$.  This is the universal property of the fibre product:
for $V\in\et X$, lifting $\bar y$ to $V\times_XY$ is the same as lifting
$f\circ\bar y$ to $V$, so the two points have the same \'etale neighbourhoods
and hence the same stalk functor.
\end{proof}

\section{Lifting \'etale morphisms along a closed immersion}
\label{sec:openclosed-lifting}

This section contains the geometric input.  It computes the stalks of $i_*$
and deduces that $i_*$ is fully faithful.

\begin{remark}[Why the dense-subsite route fails]
\label{rmk:openclosed-no-dense-subsite}
It is tempting to prove \ref{prop:openclosed-lowerStar-fully-faithful} by
showing that the base-change functor $\et X\to\et Z$ exhibits $\et Z$ as a
dense subsite, so that pushforward along it is fully faithful for formal
reasons.  The functor is indeed cover-dense --- that is exactly
\ref{lemma:openclosed-etale-lifts-locally} --- but it is neither locally full
nor locally faithful, so the dense-subsite criterion does not apply.  The
failure is real and not a matter of bookkeeping: two morphisms $V\to V'$ of
schemes \'etale over $X$ that agree after base change to $Z$ need only agree
on an open subscheme of $V$ containing the preimage of $Z$, and such an open
subscheme is not in general a covering of $V$.  The route taken below
replaces the formal argument by the stalk computation
\ref{lemma:openclosed-stalk-lowerStar}, where exactly this open subscheme is
allowed to shrink, because a stalk is a colimit over neighbourhoods.
\end{remark}

\begin{lemma}[Lifting a standard \'etale algebra along a quotient]
\label{lemma:openclosed-standard-etale-lift}
\uses{}
\lean{}
Let $R$ be a commutative ring and $I\subseteq R$ an ideal.  Recall that a
\emph{standard \'etale pair} over a ring $S$ is a pair of polynomials
$(f,g)$ in $S[T]$ with $f$ monic and $g^n\in(f',f)$ for some $n$, and that
its algebra is $S[T][1/g]/(f)$, which is \'etale over $S$.

Let $(\bar f,\bar g)$ be a standard \'etale pair over $R/I$.  Then there is a
standard \'etale pair $(f,h)$ over $R$ and an isomorphism of $R/I$-algebras
\[
  (R/I)\otimes_R\bigl(R[T][1/h]/(f)\bigr)
  \cong(R/I)[T][1/\bar g]/(\bar f).
\]
In words: every standard \'etale algebra over $R/I$ is the base change of a
standard \'etale algebra over $R$.
\end{lemma}
\begin{proof}
Lift $\bar f$ to a monic $f\in R[T]$ of the same degree, and lift the witness
of the standard-\'etale condition: choose $p_1,p_2\in R[T]$ reducing to
$\bar p_1,\bar p_2$ with
$\bar f'\bar p_1+\bar f\bar p_2=\bar g^{\,n}$, and set
\[
  h:=f'p_1+fp_2\in R[T].
\]
Then $(f,h)$ is a standard \'etale pair over $R$, with the exponent $n=1$ and
the witness $(p_1,p_2)$ available by construction; this is the whole point of
using $h$ rather than an arbitrary lift of $\bar g$, whose
standard-\'etale condition need not lift.

Base change of a standard \'etale pair along a ring map is the pair of the
reduced polynomials, so the reduction of $(f,h)$ along $R\to R/I$ is
$(\bar f,\bar g^{\,n})$.  Inverting $\bar g^{\,n}$ and inverting $\bar g$
give the same localisation, so
\[
  (R/I)[T][1/\bar g^{\,n}]/(\bar f)\cong(R/I)[T][1/\bar g]/(\bar f),
\]
which is the asserted isomorphism.
\end{proof}

\begin{lemma}[\'Etale morphisms lift locally along a closed immersion]
\label{lemma:openclosed-etale-lifts-locally}
\uses{lemma:openclosed-standard-etale-lift}
\lean{}
Let $i\colon Z\to X$ be a closed immersion, let $q\colon W\to Z$ be \'etale
and let $w\in W$.  Then there are an open subscheme $W_0\subseteq W$
containing $w$, a scheme $V$ \'etale over $X$, and an isomorphism of
$Z$-schemes
\[
  W_0\cong V\times_XZ.
\]
\end{lemma}
\begin{proof}
The assertion is local at $w$, so we may shrink $X$, $Z$ and $W$ around
$i(q(w))$ and $w$.  Choose an affine open $\Spec R\subseteq X$ containing
$i(q(w))$; then $Z$ meets it in $\Spec R/I$, and we may choose an affine open
$\Spec B\subseteq W$ containing $w$ and lying over $\Spec R/I$, so that $B$
is an \'etale $R/I$-algebra.

By the local structure of \'etale algebras, after replacing $\Spec B$ by a
basic open subscheme containing $w$ we may assume that $B$ is a standard
\'etale $R/I$-algebra.  By
\ref{lemma:openclosed-standard-etale-lift} there is a standard \'etale
$R$-algebra $A$ with $(R/I)\otimes_RA\cong B$.  Put
$V:=\Spec A$, mapping to $X$ through $\Spec A\to\Spec R\subseteq X$; this is
\'etale, being an \'etale morphism followed by an open immersion.  Finally
\[
  V\times_XZ=\Spec A\times_{\Spec R}\Spec R/I=\Spec\bigl((R/I)\otimes_RA\bigr)
  \cong\Spec B=:W_0,
\]
which is an open subscheme of $W$ containing $w$.
\end{proof}

\begin{lemma}[Stalks of $i_*$ on $Z$]
\label{lemma:openclosed-stalk-lowerStar}
\uses{lemma:openclosed-etale-lifts-locally}
\lean{}
Let $i\colon Z\to X$ be a closed immersion, let $\bar z$ be a geometric point
of $Z$ and write $\bar x:=i\circ\bar z$ for the corresponding geometric point
of $X$.  Then for every $\mathcal H\in\shv{\et Z}$,
\[
  (i_*\mathcal H)_{\bar x}\cong\mathcal H_{\bar z},
\]
naturally in $\mathcal H$.
\end{lemma}
\begin{proof}
Write $\mathrm N(\bar x)$ for the cofiltered category of \'etale
neighbourhoods $(V,\bar v)$ of $\bar x$ in $\et X$ and $\mathrm N(\bar z)$
for that of $\bar z$ in $\et Z$.  By the formula for $i_*$,
\[
  (i_*\mathcal H)_{\bar x}
  =\colim_{(V,\bar v)\in\mathrm N(\bar x)^{\mathrm{op}}}(i_*\mathcal H)(V)
  =\colim_{(V,\bar v)\in\mathrm N(\bar x)^{\mathrm{op}}}
     \mathcal H(V\times_XZ),
\]
so it suffices to prove that the base-change functor
\[
  \Phi\colon\mathrm N(\bar x)\to\mathrm N(\bar z),
  \qquad(V,\bar v)\mapsto(V\times_XZ,\bar z),
\]
is initial; the displayed isomorphism is then the induced comparison of
colimits.  By definition, initiality means that for every
$(W,\bar w)\in\mathrm N(\bar z)$ the comma category whose objects are pairs
consisting of $(V,\bar v)\in\mathrm N(\bar x)$ and a morphism
$\Phi(V,\bar v)\to(W,\bar w)$ is nonempty and connected.

\emph{Nonempty.}  This is \ref{lemma:openclosed-etale-lifts-locally} applied
to $W\to Z$ at the point underlying $\bar w$: it produces an open
neighbourhood $W_0\subseteq W$ of that point and $V$ \'etale over $X$ with
$V\times_XZ\cong W_0$.  The geometric point $\bar w$ lands in $W_0$, and its
image under $W_0\cong V\times_XZ\to V$ is a lift $\bar v$ of $\bar x$, so
$(V,\bar v)\in\mathrm N(\bar x)$; the composite
$\Phi(V,\bar v)\cong W_0\subseteq W$ is a morphism of \'etale
neighbourhoods.

\emph{Connected.}  Let $(V_1,\phi_1)$ and $(V_2,\phi_2)$ be two objects of
the comma category.  Put $V_3:=V_1\times_XV_2$, with the lift $\bar v_3$ of
$\bar x$ determined by $\bar v_1$ and $\bar v_2$; this is again an \'etale
neighbourhood of $\bar x$, and it maps to $V_1$ and to $V_2$.  Since $\bar x$
factors through $Z$, the point $\bar v_3$ also lifts to $\Phi(V_3,\bar v_3)=
V_3\times_XZ$.  The two composites
\[
  \psi_1,\psi_2\colon V_3\times_XZ\longrightarrow W
\]
obtained from $\phi_1$ and from $\phi_2$ both send that point to $\bar w$.
Now $W\to Z$ is \'etale, hence unramified, so its diagonal is an open
immersion and the locus in $V_3\times_XZ$ where two $Z$-morphisms into $W$
agree is open.  Therefore $\psi_1=\psi_2$ on an open subscheme
$\Omega\subseteq V_3\times_XZ$ containing that point.  Finally
$V_3\times_XZ\to V_3$ is a closed immersion, so $\Omega$ is the trace of an
open subscheme $\Omega'\subseteq V_3$ containing $\bar v_3$, that is,
$\Phi(\Omega',\bar v_3)=\Omega$.  The object $(\Omega',\psi_1|_\Omega)$ of the
comma category then admits morphisms to $(V_1,\phi_1)$ and to
$(V_2,\phi_2)$, given by $\Omega'\subseteq V_3\to V_i$.  Hence the comma
category is connected.

This last step is precisely where the shrinking that
\ref{rmk:openclosed-no-dense-subsite} says is unavailable to the
dense-subsite argument is harmless: in a colimit over neighbourhoods one is
allowed to pass to $\Omega'$.
\end{proof}

\begin{proposition}[$i_*$ is fully faithful]
\label{prop:openclosed-lowerStar-fully-faithful}
\uses{lemma:openclosed-stalk-lowerStar,lemma:openclosed-stalk-upperStar,lemma:openclosed-point-ulift}
\lean{}
Let $i\colon Z\to X$ be a closed immersion.  Then the counit
\[
  \varepsilon\colon i^*i_*\Longrightarrow\mathrm{id}_{\shv{\et Z}}
\]
is an isomorphism; equivalently, $i_*$ is fully faithful.
\end{proposition}
\begin{proof}
Fix $\mathcal H\in\shv{\et Z}$.  By \ref{lemma:openclosed-point-ulift} it is
enough to show that $\varepsilon_{\mathcal H}$ is an isomorphism on the stalk
at every geometric point $\bar z$ of $Z$.  By
\ref{lemma:openclosed-stalk-upperStar} applied to $i$,
\[
  (i^*i_*\mathcal H)_{\bar z}\cong(i_*\mathcal H)_{i\circ\bar z},
\]
and by \ref{lemma:openclosed-stalk-lowerStar} the right-hand side is
$\mathcal H_{\bar z}$.  Both isomorphisms are the comparison maps of the
colimit descriptions of the two sides, and the composite is the stalk of
$\varepsilon_{\mathcal H}$: unwinding the adjunction, the stalk of the counit
at $\bar z$ is the map induced on colimits by the base-change functor $\Phi$
of \ref{lemma:openclosed-stalk-lowerStar}, which is the comparison map of an
initial functor.  Hence $(\varepsilon_{\mathcal H})_{\bar z}$ is an
isomorphism for every $\bar z$, and $\varepsilon_{\mathcal H}$ is an
isomorphism.
\end{proof}

\section{Conservativity and the exact sequence}\label{sec:openclosed-exact}

\begin{proposition}[$j^*$ and $i^*$ jointly reflect isomorphisms]
\label{prop:openclosed-conservative}
\uses{lemma:openclosed-point-factors,lemma:openclosed-stalk-upperStar,lemma:openclosed-point-ulift}
\lean{}
Let $(j,i)$ be a complementary pair on $X$ and let
$\varphi\colon\mathcal F\to\mathcal F'$ be a morphism in $\shv{\et X}$ such
that $j^*\varphi$ and $i^*\varphi$ are isomorphisms.  Then $\varphi$ is an
isomorphism.
\end{proposition}
\begin{proof}
By \ref{lemma:openclosed-point-ulift} it suffices to prove that
$\varphi_{\bar x}$ is an isomorphism for every geometric point $\bar x$ of
$X$.  By \ref{lemma:openclosed-point-factors}, $\bar x$ factors as
$j\circ\bar u$ for a geometric point $\bar u$ of $U$, or as $i\circ\bar z$
for a geometric point $\bar z$ of $Z$.  In the first case
\ref{lemma:openclosed-stalk-upperStar} identifies $\varphi_{\bar x}$ with
$(j^*\varphi)_{\bar u}$, an isomorphism by hypothesis; in the second case it
identifies $\varphi_{\bar x}$ with $(i^*\varphi)_{\bar z}$, again an
isomorphism.
\end{proof}

\begin{theorem}[The open--closed exact sequence]
\label{thm:openclosed-exact}
\uses{lemma:openclosed-vanishing-lowerStar,lemma:openclosed-vanishing-lowerShriek,prop:openclosed-lowerStar-fully-faithful,prop:openclosed-conservative}
\lean{}
Let $(j,i)$ be a complementary pair on $X$ and let
$\mathcal F\in\shv{\et X}$.  Then
\[
  0\longrightarrow j_!j^*\mathcal F
   \xrightarrow{\ \alpha_{\mathcal F}\ }\mathcal F
   \xrightarrow{\ \beta_{\mathcal F}\ }i_*i^*\mathcal F\longrightarrow0
\]
is a short exact sequence in $\shv{\et X}$, naturally in $\mathcal F$.
\end{theorem}
\begin{proof}
\emph{It is a complex.}  By the adjunction $j_!\dashv j^*$ and
\ref{lemma:openclosed-vanishing-lowerStar},
\[
  \Hom(j_!j^*\mathcal F,i_*i^*\mathcal F)
  \cong\Hom(j^*\mathcal F,j^*i_*i^*\mathcal F)=0,
\]
so $\beta_{\mathcal F}\circ\alpha_{\mathcal F}=0$ with no computation.

\emph{After $j^*$.}  The functor $j^*$ is exact.  Applying it to the complex
and using that the unit $\mathrm{id}\to j^*j_!$ is an isomorphism and that
$j^*i_*=0$ by \ref{lemma:openclosed-vanishing-lowerStar}, the complex becomes
\[
  0\to j^*\mathcal F\to j^*\mathcal F\to0\to0,
\]
in which the first map is the identity by the triangle identity for
$j_!\dashv j^*$.  This is short exact.

\emph{After $i^*$.}  The functor $i^*$ is exact.  Applying it and using
$i^*j_!=0$ from \ref{lemma:openclosed-vanishing-lowerShriek} and the counit
isomorphism $i^*i_*\cong\mathrm{id}$ from
\ref{prop:openclosed-lowerStar-fully-faithful}, the complex becomes
\[
  0\to0\to i^*\mathcal F\to i^*\mathcal F\to0,
\]
in which the second map is the identity, again by a triangle identity.  This
is short exact.

\emph{Descent of exactness.}  Let $M$ be any of the three homology objects of
the complex.  Since $j^*$ and $i^*$ are exact they commute with homology, so
$j^*M=0$ and $i^*M=0$ by the two previous paragraphs.  The map $0\to M$ is
therefore an isomorphism after $j^*$ and after $i^*$, hence an isomorphism by
\ref{prop:openclosed-conservative}, that is, $M=0$.  All three homology
objects vanish, so $\alpha_{\mathcal F}$ is a monomorphism,
$\beta_{\mathcal F}$ is an epimorphism, and $\alpha_{\mathcal F}$ is a kernel
of $\beta_{\mathcal F}$.

Naturality in $\mathcal F$ holds because $\alpha$ and $\beta$ are a counit
and a unit.
\end{proof}
